\documentclass[a4paper,12pt]{article}

\usepackage[top=1cm, bottom=1.8cm, left=1cm, right=1cm]{geometry}
\usepackage[fleqn]{amsmath}
\usepackage{amssymb}
\setlength{\mathindent}{2em}
\usepackage{fontspec}
\usepackage{enumitem}
\usepackage{framed}
\usepackage{needspace}

% Set fonts
\setmainfont{Times New Roman}
\setsansfont{Arial}
\setmonofont{JetBrains Mono}

% Chinese font support
\usepackage{xeCJK}
\setCJKmainfont{Iansui}[
  BoldFont=Iansui,
  AutoFakeBold=3.17
]

\pagestyle{plain}
\newcommand{\examdate}{2026/03/06}
\newenvironment{solutionbox}{\begin{framed}}{\end{framed}}

\begin{document}

\begin{center}
{\Large\bfseries 數學試題解答}
\end{center}

\begin{flushright}
\examdate
\end{flushright}

\vspace{1em}

\begin{enumerate}[label=\textbf{\arabic*.}, leftmargin=2em, itemsep=1.2em]

\Needspace{14\baselineskip}
\item 解方程式：$2x^2 - 5x - 3 = 0$

\begin{solutionbox}
用公式
\[
x=\frac{-b\pm\sqrt{b^2-4ac}}{2a},
\]
其中 $a=2$，$b=-5$，$c=-3$。
\[
\Delta=b^2-4ac=(-5)^2-4(2)(-3)=25+24=49
\]
\[
x=\frac{-(-5)\pm\sqrt{49}}{2\cdot 2}
=\frac{5\pm 7}{4}
\]
\[
x=\frac{12}{4}=3
\quad \text{或} \quad
x=\frac{-2}{4}=-\frac{1}{2}
\]
\textbf{答案：$x=3$ 或 $x=-\dfrac{1}{2}$。}
\end{solutionbox}

\Needspace{12\baselineskip}
\item 解方程式：$x^2 - 6x + 5 = 0$

\begin{solutionbox}
因式分解：
\[
x^2-6x+5=(x-1)(x-5)=0
\]
所以
\[
x-1=0 \Rightarrow x=1,
\qquad
x-5=0 \Rightarrow x=5
\]
\textbf{答案：$x=1$ 或 $x=5$。}
\end{solutionbox}

\Needspace{10\baselineskip}
\item 將下式配方成「完全平方 $+$ 常數」：$x^2 - 8x + 3$

\begin{solutionbox}
\[
x^2-8x+3=(x^2-8x+16)-16+3
\]
\[
=(x-4)^2-13
\]
\textbf{答案：$(x-4)^2-13$。}
\end{solutionbox}

\Needspace{14\baselineskip}
\item 將下式配方成「完全平方 $+$ 常數」：$3x^2 + 12x - 7$

\begin{solutionbox}
先提出 $3$：
\[
3x^2+12x-7=3(x^2+4x)-7
\]
括號內配方：
\[
x^2+4x=(x^2+4x+4)-4=(x+2)^2-4
\]
代回得
\[
3[(x+2)^2-4]-7=3(x+2)^2-12-7=3(x+2)^2-19
\]
\textbf{答案：$3(x+2)^2-19$。}
\end{solutionbox}

\Needspace{10\baselineskip}
\item 已知點 $P(2,-1)$ 經平移向量 $(-3,5)$ 後得到點 $P'$。求 $P'$ 的座標。

\begin{solutionbox}
平移就是座標分別相加：
\[
P'=(2+(-3),-1+5)=(-1,4)
\]
\textbf{答案：$P'(-1,4)$。}
\end{solutionbox}

\Needspace{18\baselineskip}
\item 將曲線 $(x-4)^2 + (y+1)^2 = 9$ 視為某個以原點為圓心的圓經平移得到。寫出：

\begin{enumerate}[label=(\arabic*)]
\item 原本圓的方程式
\item 平移向量 $(h,k)$
\end{enumerate}

\begin{solutionbox}
由標準式
\[
(x-h)^2+(y-k)^2=r^2
\]
可知此題圓心為 $(4,-1)$，半徑平方為 $9$。

\begin{enumerate}[label=(\arabic*)]
\item 原本以原點為圓心、半徑相同的圓方程式是
\[
x^2+y^2=9
\]
\item 從 $(0,0)$ 平移到 $(4,-1)$，所以平移向量是
\[
(h,k)=(4,-1)
\]
\end{enumerate}

\textbf{答案：}
\[
x^2+y^2=9, \qquad (h,k)=(4,-1)
\]
\end{solutionbox}

\Needspace{16\baselineskip}
\item 解聯立方程式：
\[
\begin{cases}
2x + 3y = 7 \\
x - y = 1
\end{cases}
\]

\begin{solutionbox}
由第二式得
\[
x=y+1
\]
代入第一式：
\[
2(y+1)+3y=7
\Rightarrow 2y+2+3y=7
\Rightarrow 5y=5
\Rightarrow y=1
\]
\[
x=y+1=2
\]
\textbf{答案：$(x,y)=(2,1)$。}
\end{solutionbox}

\Needspace{17\baselineskip}
\item 解聯立方程式：
\[
\begin{cases}
3x - 2y = 4 \\
5x + y = 13
\end{cases}
\]

\begin{solutionbox}
由第二式得
\[
y=13-5x
\]
代入第一式：
\[
3x-2(13-5x)=4
\Rightarrow 3x-26+10x=4
\Rightarrow 13x=30
\Rightarrow x=\frac{30}{13}
\]
再求 $y$：
\[
y=13-5\cdot\frac{30}{13}
=\frac{169}{13}-\frac{150}{13}
=\frac{19}{13}
\]
\textbf{答案：}
\[
(x,y)=\left(\frac{30}{13},\frac{19}{13}\right)
\]
\end{solutionbox}

\end{enumerate}

\end{document}
